\section*{الفصل \ref{ch:b3:mantle-convection} --- الحمل في الوشاح وقوى دفع الصفائح}

\begin{solution}{exo:b3:mantle-convection:1}
$\mathrm{Ra} = 3500\times9.9\times\num{3e-5}\times1000\times(\num{6.7e5})^3/
(10^{-6}\times\num{5e20}) = \num{6.25e5}$، نحو 950 ضعف \omterm{def:b3:mantle-convection:critical}{عدد رايلي الحرج}
657.5: تحمل الطبقة حملًا قويًّا.
\end{solution}

\begin{solution}{exo:b3:mantle-convection:2}
$\tau_M = 10^{21}/\num{7.7e10} = \qty{1.30e10}{s}$، نحو \qty{412}{yr}. أمّا المدّ
الذي يدوم ساعات فيلتقي بوشاح مرن؛ وأمّا ذوبان الغطاء الجليدي
على مدى آلاف السنين، وهبوط اللوح المندس على مدى ملايين السنين، فيلتقيان
بوشاح لزج.
\end{solution}

\begin{solution}{exo:b3:mantle-convection:3}
$\tau = 4\pi\times10^{21}/(3300\times9.8\lambda)$: \qty{3.89e11}{s} =
\qty{12.3}{kyr} من أجل \qty{1000}{km}، و\qty{4.1}{kyr} من أجل \qty{3000}{km}. أمّا الحمل
الأوسع فيتعافى أسرع.
\end{solution}

\begin{solution}{exo:b3:mantle-convection:4}
$g\rho_m\alpha(T_m - T_0)\kappa t = 9.8\times3300\times\num{3.1e-5}\times1450
\times\num{0.8e-6}\times50\times\num{3.156e13} = \qty{1.835e12}{N.m^{-1}}$؛ والقوس
يساوي $1 + 2\times3300\times\num{3.1e-5}\times1450/(\pi\times2270) = 1.0416$، لذا
$F_{rp} = \qty{1.91e12}{N.m^{-1}}$، أصغر بنحو 24 مرّة من شدّ اللوح المندس.
\end{solution}

\begin{solution}{exo:b3:mantle-convection:5}
(أ) $v = 3300\times9.8\times\num{3e-5}\times300\times(\num{e4})^2/10^{21} =
\qty{2.91e-11}{m.s^{-1}}$؛ زمن الصعود عبر \qty{10}{km} هو $d/v =
\qty{3.44e14}{s}$، \qty{10.9}{Myr}، وزمن الانتشار $d^2/\kappa =
\qty{e14}{s}$، \qty{3.2}{Myr}. ونسبتهما هي $(d^2/\kappa)/(d/v) =
vd/\kappa = \rho g\alpha\Delta Td^3/(\kappa\eta) = \mathrm{Ra} = 0.29$. (ب) يبرد
أوّلًا: الكتل الصغيرة لا يحدث فيها الحمل. (ج) $\mathrm{Ra} \propto d^3$، لذا $d =
10\times(657.5/0.291)^{1/3} = \qty{131}{km}$.
\end{solution}

\begin{solution}{exo:b3:mantle-convection:6}
$2\mathrm{Ra}_c = 1315$ و$\mathrm{Ra}/2\mathrm{Ra}_c = 7604$، والتي جذرها التكعيبي
يساوي 19.66. $\delta = 2891/19.66 = \qty{147}{km}$، $\mathrm{Nu} = 19.66/2 =
9.8$, $u = (10^{-6}/\num{2.891e6})\times19.66^2 = \qty{1.34e-10}{m.s^{-1}} =
\qty{0.42}{cm/yr}$، و$q = 9.8\times3.1\times1450/\num{2.891e6} =
\qty{15}{mW.m^{-2}}$.
\end{solution}

\begin{solution}{exo:b3:mantle-convection:7}
$u = 0.09/\num{3.156e7} = \qty{2.85e-9}{m.s^{-1}}$ و$B = 2270\times\num{8e8}
\times\num{2.85e-9} = \qty{5.18e3}{kg.s^{-1}} = \qty{5.2}{Mg/s}$، نحو
ثلاثة أخماس تدفّق هاواي. ويساوي التدفّق الحراري $1170\times5178/\num{3.1e-5} =
\qty{1.95e11}{W}$، \qty{0.2}{TW}، أو \qty{0.42}{\percent} من الفقد الحراري
للأرض.
\end{solution}

\begin{solution}{exo:b3:mantle-convection:8}
$v = 2\times30\times9.8\times(\num{4e5})^2/(9\times10^{21}) =
\qty{1.05e-8}{m.s^{-1}} = \qty{0.33}{m/yr}$؛ ويستغرق قطع \qty{2000}{km}
زمنًا $\num{2e6}/\num{1.05e-8} = \qty{1.9e14}{s}$، نحو \qty{6}{Myr}.
\end{solution}

\begin{solution}{exo:b3:mantle-convection:9}
$w = 300e^{-10/4.4} = \qty{30.9}{m}$ و$\dot w = w/\tau = 30.9/4400 =
\qty{7.0}{mm/yr}$، وهو ثلثا المعدّل المقيس قرب أوميو. ونصف الفضاء المتجانس
بهذا $\tau$ يقع في النطاق الصحيح، لكنّ الارتداد الحقيقي
أسرع، أو أنّ انخفاضه أعمق.
\end{solution}

\begin{solution}{exo:b3:mantle-convection:10}
(أ) بوضع $w = 2\rho_m\alpha\Delta T\sqrt{\kappa t/\pi}/(\rho_m - \rho_w)$،
\[
  \tfrac12(\rho_m - \rho_w)gw^2 = \frac{2g\rho_m^2\alpha^2\Delta T^2\kappa t}
  {\pi(\rho_m - \rho_w)} = g\rho_m\alpha\Delta T\kappa t\times
  \frac{2\rho_m\alpha\Delta T}{\pi(\rho_m - \rho_w)} .
\] (ب) بالتكامل بالتجزئة،
$\int_0^\infty x\operatorname{erfc}x\,\mathrm dx = [\tfrac{x^2}2\operatorname{erfc}
x]_0^\infty + \int_0^\infty\tfrac{x^2}2\,\tfrac2{\sqrt\pi}e^{-x^2}\,\mathrm dx =
0 + \tfrac1{\sqrt\pi}\cdot\tfrac{\sqrt\pi}4 = \tfrac14$. (ج) عند \qty{100}{Myr}،
$\rho_mg\alpha\int z\theta = \qty{3.67e12}{N.m^{-1}}$ وبوضع $w =
\qty{3705}{m}$، $\tfrac12(\rho_m - \rho_w)gw^2 = \qty{1.53e11}{N.m^{-1}}$: حدّ
الماء لا يمثّل سوى \qty{4}{\percent} من القوّة.
\end{solution}

\begin{solution}{exo:b3:mantle-convection:11}
$\ln f = 3\ln(\pi^2 + k^2) - 2\ln k$، ومشتقّتها $6k/(\pi^2 + k^2) - 2/k$
تنعدم عندما $3k^2 = \pi^2 + k^2$، $k^2 = \pi^2/2$. وعندئذٍ $f = (3\pi^2/2)^3/
(\pi^2/2) = 27\pi^4/4$. يضمّ طول الموجة $2\pi/k = 2\sqrt2\,d = 2.83\,d$ زوجًا
من الخلايا المتعاكسة الدوران، عرض كلّ منها $\sqrt2\,d = 1.41\,d$.
\end{solution}

\begin{solution}{exo:b3:mantle-convection:12}
(أ) $u = \qty{1.90e-9}{m.s^{-1}}$، $\eta u/h = 10^{20}\times\num{1.90e-9}/
\num{1.5e5} = \qty{1.27}{MPa}$، وعلى مدى \qty{6000}{km}، \qty{7.6e12}{N.m^{-1}}.
(ب) ضعف \omterm{prop:b3:mantle-convection:ridge-push}{قوة دفع الظهرة} وسدس شدّ اللوح المندس: يستطيع \omterm{def:b3:mantle-convection:basal-drag}{السحب القاعدي} أن يقاوم
\omterm{prop:b3:mantle-convection:ridge-push}{قوة دفع الظهرة} وحدها، لا اللوح المندس. (ج) $\eta uL/h = \num{3.8e12} +
\num{4.6e12}$ يعطي $u = \num{8.4e12}\times\num{1.5e5}/(\num{6e6}\times10^{20})
= \qty{2.1e-9}{m.s^{-1}} = \qty{6.6}{cm/yr}$، وهي سرعة الصفائح المحيطية
السريعة.
\end{solution}

\begin{solution}{pb:b3:mantle-convection:1}
\textbf{1.} $p = 1000\times9.8\times300 = \qty{2.94}{MPa}$.

\textbf{2.} $w = \rho_wh/\rho_m = 300\times1000/3300 = \qty{90.9}{m}$. تهبط القشرة
ككتلة واحدة؛ وما يتدفّق من تحت الحمل ثم يعود هو
الوشاح، لذا يُوازَن العمود بالوشاح المُزاح في العمق.

\textbf{3.} $\pi\times(\num{1.25e5})^2 = \qty{4.91e10}{m^2}$، والماء
يزن $\num{4.91e10}\times300\times1000 = \qty{1.47e16}{kg}$.

\textbf{4.} تشكّل خط الشاطئ أفقيًّا، حين كان القاع أشدّ انخفاضًا عند
المركز؛ فقد ارتفع المركز أكثر من الحافّة منذ ذلك الحين، لذا يصعد الخطّ الآن
نحو المركز، بمقدار \qty{77.2}{m} على امتداد \qty{125}{km}:
\qty{0.62}{m/km}.

\textbf{5.} $\lambda/2 = \qty{250}{km}$، فإنّ $\lambda = \qty{500}{km}$ و$k =
2\pi/\num{5e5} = \qty{1.257e-5}{m^{-1}}$.

\textbf{6.} $1/k = \lambda/2\pi = \qty{79.6}{km}$.

\textbf{7.} $w = 90.9 - \text{الرفع}$: 79.4، 69.3، 52.9، 40.4، 26.9،
و\qty{13.7}{m}.

\textbf{8.} $w = w_0e^{-t/\tau}$ يعطي $\ln w = \ln w_0 - t/\tau$، خطًّا مستقيمًا
ميله $-1/\tau$.

\textbf{9.} $\tau = (14 - 1)/\ln(79.4/13.7) = 13/1.7571 = \qty{7.40}{kyr}$.

\textbf{10.} $w(6) = 79.4\,e^{-5/7.40} = \qty{40.4}{m}$، بمقدار رفع
\qty{50.5}{m}، كما في الجدول.

\textbf{11.} تبقى \qty{13.7}{m}؛ $\dot w = 13.7/7400 = \qty{1.85}{mm/yr}$.

\textbf{12.} كامل \qty{13.7}{m}، يُقترب منه ببطء متزايد.

\textbf{13.} يتناقص المعدّل وفق $e^{-t/\tau}$: $t = 7.40\ln(1.85/0.5) =
7.40\times1.309 = \qty{9.7}{kyr}$، أي نحو 9700 سنة إضافية.

\textbf{14.} $\tau = 7399\times\num{3.156e7} = \qty{2.335e11}{s}$، و$\eta =
\tau\rho_mg\lambda/(4\pi) = \num{2.335e11}\times3300\times9.8\times\num{5e5}/
12.566 = \qty{3.0e20}{Pa.s}$.

\textbf{15.} أدنى قليلًا من الطرف الأدنى لمدى فينوسكانديا. يتلاشى الجريان
أسفل حمل طول موجته \qty{500}{km} على امتداد \qty{80}{km}: يُعاين الارتداد
أضعف الوشاح وأقربه إلى السطح، أسفل ما هو في الأرض الحقيقية جزئيًّا،
داخل الغلاف الصخري.

\textbf{16.} $\tau = 4\pi\times\num{3.0e20}/(3300\times9.8\times\num{3e6}) =
\qty{3.9e10}{s} = \qty{1.2}{kyr}$. تحت حمل بهذا الاتّساع يصل الجريان
إلى عمق \qty{477}{km} فأكثر، نحو وشاح سفليّ أكثر لزوجةً بعشر مرّات،
وهو ما يُبطئه.

\textbf{17.} يحمل جزءًا من حمل ضيّق بالانثناء: الانخفاض عند
التوازن أقلّ عمقًا وأوسع من الانخفاض الأيزوستاتي البالغ \qty{90.9}{m}، وجزءٌ
من الارتداد مرنٌ وفوريّ. نصفُ فضاءٍ لزِجٍ متجانسٍ يُلائِم الرفعَ
ثم يُعطي لزوجةً تمزج التأثيرين معًا.

\textbf{18.} $\tau_M = \num{3.0e20}/\num{7.7e10} = \qty{3.9e9}{s}$، نحو
\qty{124}{yr}، أقصر بستّين مرّة من $\tau$: على مدى آلاف السنين يستجيب
الوشاح استجابةَ مائعٍ لزِج.

\textbf{19.} $1 - e^{-50/7.40} = 0.9988$: بلغ \qty{99.9}{\percent} من
الانخفاض عند التوازن، والافتراض صحيح.

\textbf{20.} $\eta \propto \lambda$: تنصيف $\lambda$ يُنصّف $\eta$،
و$\log_{10}\eta$ يتناقص بمقدار $\log_{10}2 = 0.30$.

\textbf{21.} $\mathrm{Ra} = 3700\times9.9\times\num{3.1e-5}\times1450\times
(\num{6.7e5})^3/(10^{-6}\times\num{3.0e20}) = \num{1.65e6}$.

\textbf{22.} $\num{1.65e6}/657.5 \approx 2500$ مرّة قدر الحرج.

\textbf{23.} $(\mathrm{Ra}/2\mathrm{Ra}_c)^{1/3} = (1254)^{1/3} = 10.78$: $\delta
= 670/10.78 = \qty{62}{km}$ و$\mathrm{Nu} = 5.4$.

\textbf{24.} $u = (10^{-6}/\num{6.7e5})\times10.78^2 = \qty{1.73e-10}{m.s^{-1}}
= \qty{0.55}{cm/yr}$، أبطأ من الصفائح بما بين مرّتين وأربع عشرة مرّة.

\textbf{25.} $q = 5.4\times3.1\times1450/\num{6.7e5} = \qty{36}{mW.m^{-2}}$. تُهمل
الطبقةُ الحرارةَ المتولّدةَ داخل الوشاح والحرارةَ الآتية من اللبّ والوشاح السفليّ،
والقياسُ يُسقط ثوابتَه من رتبة الواحد،
بينما تفقد الصفائح، إذ تبرد بوصفها أنصاف فضاءات منطلقةً من الظهرات، حرارةً أكبر بكثير على القاع الفتيّ
مما كانت ستفقده طبقة حدّية متوسّطة.

\textbf{26.} $\eta = \qty{3.0e20}{Pa.s}$: $\log_{10}\eta =
\textbf{\num[text-series-to-math=true]{20.48}}$.
\end{solution}
